\documentclass[12pt,a4paper]{article}

\usepackage[utf8]{inputenc}
\usepackage[spanish]{babel}
\usepackage{amsmath, amssymb, amsthm}
\usepackage{newunicodechar}
\newunicodechar{≝}{\triangleq}
\usepackage{hyperref}
\usepackage[most]{tcolorbox}
\usepackage{ocgx2}

% Definición de entornos para postulados y teoremas con tcolorbox
\tcbuselibrary{theorems}

\newtcbtheorem{preaxioma}{Pre-Axioma}{
  colback=purple!5!white,
  colframe=purple!60!black,
  fonttitle=\bfseries,
  arc=2mm,
  boxrule=1pt
}{prax}

\newtcbtheorem{postulado}{Postulado}{
  colback=blue!5!white,
  colframe=blue!60!black,
  fonttitle=\bfseries,
  arc=2mm,
  boxrule=1pt
}{post}

\newtcbtheorem{teorema}{Teorema}{
  colback=green!5!white,
  colframe=green!50!black,
  fonttitle=\bfseries,
  arc=2mm,
  boxrule=1pt
}{teo}

\newtcbtheorem{definicion}{Definición}{
  colback=cyan!5!white,
  colframe=cyan!60!black,
  fonttitle=\bfseries,
  arc=2mm,
  boxrule=1pt
}{def}

% Macro para pruebas desplegables
\newcommand{\desplegable}[2]{%
  \par\vspace{1ex}\noindent
  \switchocg{#1}{\fcolorbox{gray!50}{gray!10}{\textsf{\textbf{\small $\blacktriangledown$ Mostrar/Ocultar Demostración}}}}
  \begin{ocg}{Demostración}{#1}{off}
  \vspace{1ex}
  #2
  \end{ocg}
  \par\vspace{1ex}
}

\title{Desarrollo del Álgebra de Boole \\ \large Desde los Postulados de Huntington}
\author{}
\date{\today}

\begin{document}

\maketitle

\section{Introducción}
El álgebra de Boole, originalmente introducida por George Boole en su obra de 1854 \textit{The Laws of Thought}, fue posteriormente axiomatizada de forma rigurosa. En 1904, el matemático estadounidense Edward Vermilye Huntington presentó un conjunto de postulados independientes que definen formalmente una estructura de álgebra de Boole. 

Para evitar cualquier confusión conceptual con la aritmética tradicional, en esta fase inicial emplearemos la signatura propia de la teoría de retículos, utilizando los símbolos $\vee$ (supremo o join) y $\wedge$ (ínfimo o meet), junto a los elementos constantes $\bot$ (mínimo) y $\top$ (máximo). El complemento se denotará con el símbolo clásico de la negación lógica $\neg$. Más adelante, y por conveniencia práctica, transitaremos hacia la notación clásica de sistemas digitales ($+$, $\cdot$, $0$, $1$).

\section{Pre-Axiomas de la Estructura}
Antes de enunciar los postulados, debemos definir rigurosamente sobre qué elementos y operaciones estamos trabajando. 

Partimos de un ente matemático $\mathbb{B}$.

\begin{preaxioma}{Estructura de Conjunto - EsConjunto($\mathbb{B}$)}{esconj}
Se requiere que $\mathbb{B}$ sea un conjunto.
\end{preaxioma}

\begin{preaxioma}{Elementos Constantes - Constantes}{constantes}
Este conjunto ha de cumplir que tiene dos elementos que llamaremos constantes, tales que $\bot \in \mathbb{B}$ y $\top \in \mathbb{B}$. En principio, no asumimos nada sobre la igualdad o desigualdad de estas constantes.
\end{preaxioma}

Además, vamos a definir dos operaciones binarias internas que denotaremos por $\vee$ y $\wedge$. Estas deben satisfacer rigurosamente la definición de función:

\begin{preaxioma}{Operación Binaria Interna $\vee$ - OpBinInt$_\vee$}{opbinint_vee}
$\vee : \mathbb{B} \times \mathbb{B} \to \mathbb{B}$ es una operación binaria interna.
\end{preaxioma}

\begin{preaxioma}{Operación Binaria Interna $\wedge$ - OpBinInt$_\wedge$}{opbinint_wedge}
$\wedge : \mathbb{B} \times \mathbb{B} \to \mathbb{B}$ es una operación binaria interna.
\end{preaxioma}

Para poder usar estos conceptos con mayor seguridad y flexibilidad en las futuras demostraciones formales, asignaremos nombres cortos a las condiciones de existencia y unicidad de la imagen para estas operaciones:

\begin{preaxioma}{Existencia $\vee$ - Existencia$_\vee$}{exist_vee}
Para todo par existe imagen en $\mathbb{B}$. Es decir, $\forall \langle a,b \rangle \in \mathbb{B} \times \mathbb{B}$, $\exists c \in \mathbb{B}$ tal que $a \vee b = c$.
\end{preaxioma}

\begin{preaxioma}{Existencia $\wedge$ - Existencia$_\wedge$}{exist_wedge}
Análogamente, $\forall \langle a,b \rangle \in \mathbb{B} \times \mathbb{B}$, $\exists d \in \mathbb{B}$ tal que $a \wedge b = d$.
\end{preaxioma}

\begin{preaxioma}{Unicidad $\vee$ - Unicidad$_\vee$}{unic_vee}
Para un par solo existe una imagen. Esto es, si $a \vee b = c$ y $a \vee b = d$, entonces $c = d$.
\end{preaxioma}

\begin{preaxioma}{Unicidad $\wedge$ - Unicidad$_\wedge$}{unic_wedge}
De igual forma para el ínfimo, si $a \wedge b = c$ y $a \wedge b = d$, entonces $c = d$.
\end{preaxioma}

\section{Los Postulados de Huntington (1904)}
Sobre el sistema $(\mathbb{B}, \vee, \wedge, \bot, \top)$ que cumple los pre-axiomas anteriores, diremos que forma un álgebra de Boole si satisface los siguientes postulados:

\begin{postulado}{Elemento neutro $\vee$ - $ElemNeu_\vee$}{neutro_vee}
Todo elemento operado mediante $\vee$ con el mínimo $\bot$ da como resultado el mismo elemento; es decir, $\bot$ no altera el valor original:
\[ \forall a \in \mathbb{B}, \quad a \vee \bot = a \]
\end{postulado}

\begin{postulado}{Elemento neutro $\wedge$ - $ElemNeu_\wedge$}{neutro_wedge}
Todo elemento operado mediante $\wedge$ con el máximo $\top$ da como resultado el mismo elemento, quedando inalterado:
\[ \forall a \in \mathbb{B}, \quad a \wedge \top = a \]
\end{postulado}

\begin{postulado}{Conmutatividad $\vee$ - $Comm_\vee$}{conmut_vee}
El orden de los operandos al aplicar la operación $\vee$ es indiferente, obteniéndose exactamente el mismo resultado:
\[ \forall a, b \in \mathbb{B}, \quad a \vee b = b \vee a \]
\end{postulado}

\begin{postulado}{Conmutatividad $\wedge$ - $Comm_\wedge$}{conmut_wedge}
De la misma forma, el orden de los operandos al aplicar la operación $\wedge$ tampoco altera el resultado final:
\[ \forall a, b \in \mathbb{B}, \quad a \wedge b = b \wedge a \]
\end{postulado}

\begin{postulado}{Distributividad $\vee$ sobre $\wedge$ - $Dist_\vee$}{distrib_vee_wedge}
La operación $\vee$ se distribuye sobre la operación $\wedge$. Operar un elemento con el resultado de un $\wedge$ equivale a operar con $\vee$ cada componente individualmente y luego aplicar $\wedge$:
\[ \forall a, b, c \in \mathbb{B}, \quad a \vee (b \wedge c) = (a \vee b) \wedge (a \vee c) \]
\end{postulado}

\begin{postulado}{Distributividad $\wedge$ sobre $\vee$ - $Dist_\wedge$}{distrib_wedge_vee}
De manera equivalente, el ínfimo ($\wedge$) se reparte de forma distributiva entre los componentes de un supremo ($\vee$):
\[ \forall a, b, c \in \mathbb{B}, \quad a \wedge (b \vee c) = (a \wedge b) \vee (a \wedge c) \]
\end{postulado}

\begin{postulado}{Complementario - $Comp_\vee, Comp_\wedge$}{comp}
Todo elemento del conjunto posee al menos un "complemento" (o elemento opuesto). Al operarlo con su complemento mediante $\vee$ siempre alcanzamos el máximo $\top$, y mediante $\wedge$ siempre caemos al mínimo $\bot$:
\begin{align*}
\forall a \in \mathbb{B}, \exists b \in \mathbb{B} \quad : \quad a \vee b &= \top \quad (Comp_\vee) \\
a \wedge b &= \bot \quad (Comp_\wedge)
\end{align*}
\end{postulado}

\textit{Nota: A diferencia de algunas formulaciones clásicas que imponen un axioma de cardinalidad ($\bot \neq \top$) para evitar el álgebra trivial, en este desarrollo permitiremos la existencia del álgebra trivial.}

\section{El Principio de Dualidad}
Si observamos los postulados de Huntington, notaremos una perfecta simetría entre las operaciones $\vee$ y $\wedge$, y entre las constantes $\bot$ y $\top$. Si en cualquier postulado intercambiamos $\vee$ por $\wedge$ y $\bot$ por $\top$, obtenemos otro postulado válido del sistema. 

Este rasgo estructural da lugar al \textbf{Principio de Dualidad}: toda proposición o teorema deducido a partir de estos axiomas tiene un \textit{teorema dual} que también es válido. La demostración de un teorema dual se construye manipulando la prueba original y aplicando sistemáticamente el intercambio de operaciones ($\vee \leftrightarrow \wedge$) y constantes ($\bot \leftrightarrow \top$). En las siguientes pruebas no evitaremos repetir las versiones duales; al contrario, haremos hincapié en cómo se manipula la prueba de uno para obtener la del otro.

\section{Teoremas Principales Derivados}

\begin{teorema}{Unicidad de los elementos neutros - $Unic_e, Unic_u$}{unicidad_neutros}
Los elementos neutros descritos en los postulados son únicos. No existe ningún otro elemento en el conjunto que se comporte como el mínimo $\bot$ para la operación $\vee$, ni ningún otro que actúe como el máximo $\top$ para la operación $\wedge$:
\begin{align*}
    \exists! e \in \mathbb{B}, \forall a \in \mathbb{B}, a \vee e &= a \implies e = \bot \quad (Unic_e) \\
    \exists! u \in \mathbb{B}, \forall a \in \mathbb{B}, a \wedge u &= a \implies u = \top \quad (Unic_u)
\end{align*}
\end{teorema}
\desplegable{proof_unic}{
\begin{proof}[Demostración de $Unic_e$]
Sea $e \in \mathbb{B}$ tal que $\forall a \in \mathbb{B}, a \vee e = a$. Tomando $a = \bot$:
\begin{align*}
    e &= e \vee \bot & (ElemNeu_\vee) \\
      &= \bot \vee e & (Comm_\vee) \\
      &= \bot & (\text{Hipótesis sobre } e)
\end{align*}
\end{proof}
\begin{proof}[Demostración de $Unic_u$ (Dual)]
Para obtener la prueba dual, intercambiamos $\vee$ por $\wedge$ y $\bot$ por $\top$.
Sea $u \in \mathbb{B}$ tal que $\forall a \in \mathbb{B}, a \wedge u = a$. Tomando $a = \top$:
\begin{align*}
    u &= u \wedge \top & (ElemNeu_\wedge) \\
      &= \top \wedge u & (Comm_\wedge) \\
      &= \top & (\text{Hipótesis sobre } u)
\end{align*}
\end{proof}
}

\begin{teorema}{Idempotencia - $Idemp_\vee, Idemp_\wedge$}{idempotencia}
Operar un elemento consigo mismo, independientemente de si usamos $\vee$ o $\wedge$, no altera su valor. El elemento se mantiene idéntico a sí mismo:
\begin{align*}
    \forall a \in \mathbb{B}, \quad a \vee a &= a \quad (Idemp_\vee) \\
    \forall a \in \mathbb{B}, \quad a \wedge a &= a \quad (Idemp_\wedge)
\end{align*}
\end{teorema}
\desplegable{proof_idemp}{
\begin{proof}[Demostración de $Idemp_\vee$]
\begin{align*}
    a &= a \vee \bot & (ElemNeu_\vee) \\
      &= a \vee (a \wedge \neg a) & (Comp_\wedge) \\
      &= (a \vee a) \wedge (a \vee \neg a) & (Dist_\vee) \\
      &= (a \vee a) \wedge \top & (Comp_\vee) \\
      &= \top \wedge (a \vee a) & (Comm_\wedge) \\
      &= a \vee a & (ElemNeu_\wedge)
\end{align*}
\end{proof}
\begin{proof}[Demostración de $Idemp_\wedge$ (Dual)]
Intercambiando los operadores y constantes de la prueba anterior paso a paso:
\begin{align*}
    a &= a \wedge \top & (ElemNeu_\wedge) \\
      &= a \wedge (a \vee \neg a) & (Comp_\vee) \\
      &= (a \wedge a) \vee (a \wedge \neg a) & (Dist_\wedge) \\
      &= (a \wedge a) \vee \bot & (Comp_\wedge) \\
      &= \bot \vee (a \wedge a) & (Comm_\vee) \\
      &= a \wedge a & (ElemNeu_\vee)
\end{align*}
\end{proof}
}

\begin{teorema}{Elementos absorbentes - $Abs_{\bot}, Abs_{\top}$}{absorbentes}
Cualquier elemento operado mediante $\vee$ con el máximo $\top$ es absorbido por este, dando como resultado $\top$. De igual manera, operar cualquier elemento mediante $\wedge$ con el mínimo $\bot$ siempre resulta en $\bot$:
\begin{align*}
    \forall a \in \mathbb{B}, \quad a \vee \top &= \top \\
    \forall a \in \mathbb{B}, \quad a \wedge \bot &= \bot
\end{align*}
\end{teorema}
\desplegable{proof_absorbentes}{
\begin{proof}[Demostración de $a \vee \top = \top$]
\begin{align*}
    a \vee \top &= (a \vee \top) \wedge \top & (ElemNeu_\wedge) \\
               &= (a \vee \top) \wedge (a \vee \neg a) & (Comp_\vee) \\
               &= a \vee (\top \wedge \neg a) & (Dist_\vee) \\
               &= a \vee (\neg a \wedge \top) & (Comm_\wedge) \\
               &= a \vee \neg a & (ElemNeu_\wedge) \\
               &= \top & (Comp_\vee)
\end{align*}
\end{proof}
\begin{proof}[Demostración de $a \wedge \bot = \bot$ (Dual)]
\begin{align*}
    a \wedge \bot &= (a \wedge \bot) \vee \bot & (ElemNeu_\vee) \\
               &= (a \wedge \bot) \vee (a \wedge \neg a) & (Comp_\wedge) \\
               &= a \wedge (\bot \vee \neg a) & (Dist_\wedge) \\
               &= a \wedge (\neg a \vee \bot) & (Comm_\vee) \\
               &= a \wedge \neg a & (ElemNeu_\vee) \\
               &= \bot & (Comp_\wedge)
\end{align*}
\end{proof}
}

\begin{teorema}{Condición de Álgebra Trivial}{trivial_cond}
Si se da el caso extremo de que el elemento mínimo $\bot$ y el máximo $\top$ son exactamente el mismo, entonces estamos ante un álgebra que contiene un único elemento en todo su conjunto (el álgebra trivial):
\[ \bot = \top \implies \mathbb{B} = \{\top\} = \{\bot\} \]
\end{teorema}
\desplegable{proof_trivial_cond}{
\begin{proof}
Supongamos que $\bot = \top$. Sea $x \in \mathbb{B}$ un elemento cualquiera:
\begin{align*}
    x &= x \vee \bot & (ElemNeu_\vee) \\
      &= x \vee \top & (\text{Hipótesis } \bot = \top) \\
      &= \top & (Abs_\top)
\end{align*}
Por tanto, todo elemento $x$ del conjunto es idéntico a $\top$, lo que implica que $\mathbb{B} = \{\top\} = \{\bot\}$.
\end{proof}
}

\begin{teorema}{Complemento Idéntico implica Álgebra Trivial}{trivial_comp}
Si dentro de la estructura existe algún elemento que sea igual a su propio complemento ($\neg a = a$), entonces forzosamente todo el sistema colapsa en el álgebra trivial de un solo elemento:
\[ (\exists a \in \mathbb{B} : \neg a = a) \implies \mathbb{B} = \{\top\} = \{\bot\} \]
\end{teorema}
\desplegable{proof_trivial_comp}{
\begin{proof}
Supongamos que existe $a \in \mathbb{B}$ tal que $\neg a = a$. Por el postulado del Complemento ($Comp_\vee$ y $Comp_\wedge$), sabemos que $a \vee \neg a = \top$ y $a \wedge \neg a = \bot$. Sustituyendo la hipótesis $\neg a = a$ en ambas ecuaciones, obtenemos:
$$ a \vee a = \top \quad \text{y} \quad a \wedge a = \bot $$
Aplicando el teorema de Idempotencia ($Idemp_\vee$ y $Idemp_\wedge$), sabemos que $a \vee a = a$ y $a \wedge a = a$. Por tanto:
$$ a = \top \quad \text{y} \quad a = \bot $$
Lo cual implica que $\bot = \top$. Aplicando el teorema anterior (Condición de Álgebra Trivial), concluimos que $\mathbb{B} = \{\top\} = \{\bot\}$.
\end{proof}
}

\begin{teorema}{Propiedades de absorción - $Abs_{\vee}, Abs_{\wedge}$}{absorcion}
Cuando se combinan ambas operaciones anidando un elemento consigo mismo y con un tercero, el elemento repetido "absorbe" al otro, independientemente del valor del segundo:
\begin{align*}
    \forall a, b \in \mathbb{B}, \quad a \vee (a \wedge b) &= a \\
    \forall a, b \in \mathbb{B}, \quad a \wedge (a \vee b) &= a
\end{align*}
\end{teorema}
\desplegable{proof_absorcion}{
\begin{proof}[Demostración de $Abs_\vee$]
\begin{align*}
    a \vee (a \wedge b) &= (a \wedge \top) \vee (a \wedge b) & (ElemNeu_\wedge) \\
                        &= a \wedge (\top \vee b) & (Dist_\wedge) \\
                        &= a \wedge (b \vee \top) & (Comm_\vee) \\
                        &= a \wedge \top & (Abs_\top) \\
                        &= a & (ElemNeu_\wedge)
\end{align*}
\end{proof}
\begin{proof}[Demostración de $Abs_\wedge$ (Dual)]
\begin{align*}
    a \wedge (a \vee b) &= (a \vee \bot) \wedge (a \vee b) & (ElemNeu_\vee) \\
                        &= a \vee (\bot \wedge b) & (Dist_\vee) \\
                        &= a \vee (b \wedge \bot) & (Comm_\wedge) \\
                        &= a \vee \bot & (Abs_\bot) \\
                        &= a & (ElemNeu_\vee)
\end{align*}
\end{proof}
}

\begin{teorema}{Propiedades de orden de retículo - $Prop_{\vee,\wedge}$}{prop_reticulo}
Existe una correspondencia biunívoca fundamental entre las dos operaciones: afirmar que un elemento domina a otro mediante $\vee$ equivale matemáticamente a afirmar que el segundo se impone al primero mediante $\wedge$:
\[ \forall a, b \in \mathbb{B}: \quad a \vee b = a \iff a \wedge b = b \]
\end{teorema}
\desplegable{proof_prop_reticulo}{
\begin{proof}[Demostración de $\implies$]
Supongamos que $a \vee b = a$.
\begin{align*}
    a \wedge b &= b \wedge a & (Comm_\wedge) \\
               &= b \wedge (a \vee b) & (\text{Hipótesis } a \vee b = a) \\
               &= b & (Abs_\wedge)
\end{align*}
\end{proof}
\begin{proof}[Demostración de $\impliedby$ (Dual)]
Supongamos que $a \wedge b = b$.
\begin{align*}
    a \vee b &= b \vee a & (Comm_\vee) \\
               &= (a \wedge b) \vee a & (\text{Hipótesis } a \wedge b = b) \\
               &= a \vee (a \wedge b) & (Comm_\vee) \\
               &= a & (Abs_\vee)
\end{align*}
\end{proof}
}

\begin{teorema}{Equivalencia de operaciones - $Equa_{\vee,\wedge}$}{equa_operaciones}
Si operar dos elementos mediante $\vee$ da exactamente el mismo resultado que operarlos mediante $\wedge$, esto sólo es lógicamente posible si ambos elementos son en realidad el mismo:
\[ \forall a, b \in \mathbb{B}: \quad a \vee b = a \wedge b \implies a = b \]
\end{teorema}
\desplegable{proof_equa_operaciones}{
\begin{proof}
Supongamos $a \vee b = a \wedge b$. Observamos que:
\begin{align*}
    a &= a \vee (a \wedge b) & (Abs_\vee) \\
      &= a \vee (a \vee b) & (\text{Hipótesis})
\end{align*}
Aplicando $Prop_{\vee,\wedge}$, dado que $a \vee (a \vee b) = a$, deducimos que $a \wedge (a \vee b) = a \vee b$.
\begin{align*}
    a \vee b &= a \wedge (a \vee b) & (\text{Resultado anterior}) \\
             &= a & (Abs_\wedge)
\end{align*}
De manera simétrica para $b$:
\begin{align*}
    b &= b \vee (a \wedge b) & (Abs_\vee) \\
      &= b \vee (a \vee b) & (\text{Hipótesis}) \\
      &= (a \vee b) \vee b & (Comm_\vee)
\end{align*}
Aplicando de nuevo $Prop_{\vee,\wedge}$ sobre esta igualdad, obtenemos que $(a \vee b) \wedge b = a \vee b$. Pero sabemos por $Abs_\wedge$ que $(a \vee b) \wedge b = b \wedge (a \vee b) = b$. 
Por consiguiente, $a \vee b = b$. Finalmente, uniendo ambos resultados: $a = a \vee b = b$.
\end{proof}
}

\begin{teorema}{Teorema de Cancelación - $Equa_{canc}$}{equa_canc}
Si un elemento $a$ se opera mediante $\vee$ con $b$ y con $c$ dando el mismo resultado, y además se opera mediante $\wedge$ con $b$ y con $c$ coincidiendo también los resultados, entonces forzosamente $b$ y $c$ son el mismo elemento:
\[ \forall a, b, c \in \mathbb{B}: \quad a \vee b = a \vee c \quad \text{y} \quad a \wedge b = a \wedge c \implies b = c \]
\end{teorema}
\desplegable{proof_equa_canc}{
\begin{proof}
Este teorema es automejor dualizable al ser sus hipótesis perfectamente simétricas.
\begin{align*}
    b &= b \wedge (a \vee b) & (Abs_\wedge) \\
      &= b \wedge (a \vee c) & (\text{Hipótesis } a \vee b = a \vee c) \\
      &= (b \wedge a) \vee (b \wedge c) & (Dist_\wedge) \\
      &= (a \wedge b) \vee (b \wedge c) & (Comm_\wedge) \\
      &= (a \wedge c) \vee (b \wedge c) & (\text{Hipótesis } a \wedge b = a \wedge c) \\
      &= (c \wedge a) \vee (c \wedge b) & (Comm_\wedge \text{ aplicado dos veces}) \\
      &= c \wedge (a \vee b) & (Dist_\wedge) \\
      &= c \wedge (a \vee c) & (\text{Hipótesis } a \vee b = a \vee c) \\
      &= c & (Abs_\wedge)
\end{align*}
\end{proof}
}

\begin{teorema}{Unicidad del complemento - $Unic_{comp}$}{unic_comp}
Todo elemento del conjunto tiene un complemento $\neg a$, y este es estrictamente único. Ningún otro elemento puede cumplir simultáneamente las dos condiciones del postulado del complemento para un mismo $a$:
\[ \forall a, x \in \mathbb{B} : \quad (a \vee x = \top \quad \text{y} \quad a \wedge x = \bot) \implies x = \neg a \]
\end{teorema}
\desplegable{proof_unic_comp}{
\begin{proof}
Supongamos que existe $x \in \mathbb{B}$ tal que $a \vee x = \top$ y $a \wedge x = \bot$.
\begin{align*}
    x &= x \wedge \top & (ElemNeu_\wedge) \\
      &= x \wedge (a \vee \neg a) & (Comp_\vee) \\
      &= (x \wedge a) \vee (x \wedge \neg a) & (Dist_\wedge) \\
      &= (a \wedge x) \vee (x \wedge \neg a) & (Comm_\wedge) \\
      &= \bot \vee (x \wedge \neg a) & (\text{Hipótesis } a \wedge x = \bot) \\
      &= (a \wedge \neg a) \vee (x \wedge \neg a) & (Comp_\wedge) \\
      &= (a \vee x) \wedge \neg a & (Dist_\wedge) \\
      &= \top \wedge \neg a & (\text{Hipótesis } a \vee x = \top) \\
      &= \neg a \wedge \top & (Comm_\wedge) \\
      &= \neg a & (ElemNeu_\wedge)
\end{align*}
Por tanto, si $x$ cumple las condiciones de complemento, $x$ tiene que ser necesariamente $\neg a$.
\end{proof}
}

\begin{teorema}{Involución - $Comp_{inv}$}{comp_inv}
Aplicar la operación de complemento (o negación) dos veces consecutivas sobre un mismo elemento cancela su efecto, devolviendo el elemento original intacto:
\[ \forall a \in \mathbb{B}, \quad \neg (\neg a) = a \]
\end{teorema}
\desplegable{proof_comp_inv}{
\begin{proof}
Por definición, el complemento de $\neg a$, denotado como $\neg (\neg a)$, es el elemento único que satisface:
$$ \neg a \vee \neg (\neg a) = \top \quad \text{y} \quad \neg a \wedge \neg (\neg a) = \bot $$
Sin embargo, por la conmutatividad ($Comm_\vee$ y $Comm_\wedge$), sabemos que:
$$ \neg a \vee a = a \vee \neg a = \top \quad \text{y} \quad \neg a \wedge a = a \wedge \neg a = \bot $$
Esto demuestra que $a$ actúa como un complemento de $\neg a$. Por el teorema de unicidad del complemento ($Unic_{comp}$), concluimos necesariamente que $\neg (\neg a) = a$.
\end{proof}
}

\begin{teorema}{Leyes de De Morgan - $Mor_{\vee, \wedge}$}{morgan}
La negación matemática se distribuye sobre las operaciones, pero al hacerlo, invierte la operación original: un supremo ($\vee$) negado se convierte en el ínfimo ($\wedge$) de las negaciones, y viceversa:
\begin{align}
    \neg (a \vee b) &= \neg a \wedge \neg b \\
    \neg (a \wedge b) &= \neg a \vee \neg b 
\end{align}
De forma equivalente, aislando las variables mediante la involución, podemos expresar las operaciones básicas exclusivamente a partir de su dual negada:
\begin{align}
    a \vee b &= \neg (\neg a \wedge \neg b) \\
    a \wedge b &= \neg (\neg a \vee \neg b) 
\end{align}
\end{teorema}
\desplegable{proof_morgan}{
\begin{proof}[Demostración de $\neg (a \vee b) = \neg a \wedge \neg b$]
Para demostrarlo sin recurrir a la asociatividad, usaremos las propiedades de absorción. Comprobemos primero la suma: $(a \vee b) \vee (\neg a \wedge \neg b) = \top$.
Sabemos por $Abs_\wedge$ que $a \wedge (a \vee b) = a$.
\begin{align*}
    \neg a \vee (a \wedge (a \vee b)) &= \neg a \vee a = \top & (Comp_\vee) \\
    (\neg a \vee a) \wedge (\neg a \vee (a \vee b)) &= \top & (Dist_\vee) \\
    \top \wedge (\neg a \vee (a \vee b)) &= \top & (Comp_\vee) \\
    \neg a \vee (a \vee b) &= \top & (ElemNeu_\wedge)
\end{align*}
Simétricamente, como $b \wedge (a \vee b) = b \wedge (b \vee a) = b$, obtenemos $\neg b \vee (a \vee b) = \top$.
Por tanto:
\begin{align*}
    (a \vee b) \vee (\neg a \wedge \neg b) &= ((a \vee b) \vee \neg a) \wedge ((a \vee b) \vee \neg b) & (Dist_\vee) \\
    &= (\neg a \vee (a \vee b)) \wedge (\neg b \vee (a \vee b)) & (Comm_\vee) \\
    &= \top \wedge \top & (\text{Resultados anteriores}) \\
    &= \top & (Idemp_\wedge)
\end{align*}

Segundo, comprobemos el producto: $(a \vee b) \wedge (\neg a \wedge \neg b) = \bot$.
Sabemos por $Abs_\vee$ que $\neg a \vee (\neg a \wedge \neg b) = \neg a$.
\begin{align*}
    a \wedge (\neg a \vee (\neg a \wedge \neg b)) &= a \wedge \neg a = \bot & (Comp_\wedge) \\
    (a \wedge \neg a) \vee (a \wedge (\neg a \wedge \neg b)) &= \bot & (Dist_\wedge) \\
    \bot \vee (a \wedge (\neg a \wedge \neg b)) &= \bot & (Comp_\wedge) \\
    a \wedge (\neg a \wedge \neg b) &= \bot & (ElemNeu_\vee)
\end{align*}
Simétricamente, como $\neg b \vee (\neg a \wedge \neg b) = \neg b \vee (\neg b \wedge \neg a) = \neg b$, obtenemos $b \wedge (\neg a \wedge \neg b) = \bot$.
Por tanto:
\begin{align*}
    (a \vee b) \wedge (\neg a \wedge \neg b) &= (\neg a \wedge \neg b) \wedge (a \vee b) & (Comm_\wedge) \\
    &= ((\neg a \wedge \neg b) \wedge a) \vee ((\neg a \wedge \neg b) \wedge b) & (Dist_\wedge) \\
    &= (a \wedge (\neg a \wedge \neg b)) \vee (b \wedge (\neg a \wedge \neg b)) & (Comm_\wedge) \\
    &= \bot \vee \bot & (\text{Resultados anteriores}) \\
    &= \bot & (Idemp_\vee)
\end{align*}
Por el teorema de unicidad ($Unic_{comp}$), concluimos que $\neg (a \vee b) = \neg a \wedge \neg b$.
\end{proof}

\begin{proof}[Demostración de $\neg (a \wedge b) = \neg a \vee \neg b$ (Dual)]
Intercambiando operaciones y constantes, comprobamos el producto: $(a \wedge b) \wedge (\neg a \vee \neg b) = \bot$.
Sabemos por $Abs_\vee$ que $a \vee (a \wedge b) = a$.
\begin{align*}
    \neg a \wedge (a \vee (a \wedge b)) &= \neg a \wedge a = \bot & (Comp_\wedge) \\
    (\neg a \wedge a) \vee (\neg a \wedge (a \wedge b)) &= \bot & (Dist_\wedge) \\
    \bot \vee (\neg a \wedge (a \wedge b)) &= \bot & (Comp_\wedge) \\
    \neg a \wedge (a \wedge b) &= \bot & (ElemNeu_\vee)
\end{align*}
Simétricamente, $\neg b \wedge (a \wedge b) = \bot$.
Por tanto:
\begin{align*}
    (a \wedge b) \wedge (\neg a \vee \neg b) &= ((a \wedge b) \wedge \neg a) \vee ((a \wedge b) \wedge \neg b) & (Dist_\wedge) \\
    &= (\neg a \wedge (a \wedge b)) \vee (\neg b \wedge (a \wedge b)) & (Comm_\wedge) \\
    &= \bot \vee \bot = \bot & (\text{Resultados anteriores})
\end{align*}

Comprobemos la suma: $(a \wedge b) \vee (\neg a \vee \neg b) = \top$.
Sabemos por $Abs_\wedge$ que $\neg a \wedge (\neg a \vee \neg b) = \neg a$.
\begin{align*}
    a \vee (\neg a \wedge (\neg a \vee \neg b)) &= a \vee \neg a = \top & (Comp_\vee) \\
    (a \vee \neg a) \wedge (a \vee (\neg a \vee \neg b)) &= \top & (Dist_\vee) \\
    \top \wedge (a \vee (\neg a \vee \neg b)) &= \top & (Comp_\vee) \\
    a \vee (\neg a \vee \neg b) &= \top & (ElemNeu_\wedge)
\end{align*}
Simétricamente, $b \vee (\neg a \vee \neg b) = \top$.
Por tanto:
\begin{align*}
    (a \wedge b) \vee (\neg a \vee \neg b) &= (\neg a \vee \neg b) \vee (a \wedge b) & (Comm_\vee) \\
    &= ((\neg a \vee \neg b) \vee a) \wedge ((\neg a \vee \neg b) \vee b) & (Dist_\vee) \\
    &= (a \vee (\neg a \vee \neg b)) \wedge (b \vee (\neg a \vee \neg b)) & (Comm_\vee) \\
    &= \top \wedge \top = \top & (\text{Resultados anteriores})
\end{align*}
Por $Unic_{comp}$, $\neg (a \wedge b) = \neg a \vee \neg b$.
\end{proof}

\begin{proof}[Demostración de $a \vee b = \neg (\neg a \wedge \neg b)$ y su dual]
Partiendo de $\neg (\neg a \wedge \neg b)$, aplicamos De Morgan a sus componentes:
\begin{align*}
    \neg (\neg a \wedge \neg b) &= \neg (\neg a) \vee \neg (\neg b) & (Mor_\wedge) \\
    &= a \vee b & (Comp_{inv})
\end{align*}
Dualizando la expresión, obtenemos de manera idéntica que $\neg (\neg a \vee \neg b) = a \wedge b$.
\end{proof}
}

\begin{teorema}{Asociatividad - $Asoc_\vee, Asoc_\wedge$}{asociatividad}
El orden en el que se agrupan tres o más elementos al aplicar de forma consecutiva la misma operación ($\vee$ o $\wedge$) no altera el resultado final. Al ubicar este teorema después de De Morgan, podemos simplificar enormemente su demostración:
\begin{align*}
    a \vee (b \vee c) &= (a \vee b) \vee c \\
    a \wedge (b \wedge c) &= (a \wedge b) \wedge c
\end{align*}
\end{teorema}
\desplegable{proof_asociatividad}{
\begin{proof}[Demostración de $a \vee (b \vee c) = (a \vee b) \vee c$]
Primero, demostraremos un pequeño \textbf{Lema de Igualdad por Casos}: Si $x \wedge y = x \wedge z$ y $\neg x \wedge y = \neg x \wedge z$, entonces $y = z$.
\begin{align*}
    y &= \top \wedge y & (ElemNeu_\wedge) \\
      &= (x \vee \neg x) \wedge y & (Comp_\vee) \\
      &= (x \wedge y) \vee (\neg x \wedge y) & (Dist_\wedge) \\
      &= (x \wedge z) \vee (\neg x \wedge z) & (\text{Por hipótesis del Lema}) \\
      &= (x \vee \neg x) \wedge z & (Dist_\wedge) \\
      &= \top \wedge z = z & (Comp_\vee, ElemNeu_\wedge)
\end{align*}
Sea $L = a \vee (b \vee c)$ y $R = (a \vee b) \vee c$. Aplicaremos el lema usando $x = a$, por lo que debemos demostrar que $a \wedge L = a \wedge R$ y $\neg a \wedge L = \neg a \wedge R$.

1) Comprobamos $a \wedge L = a \wedge R$:
\begin{align*}
    a \wedge L &= a \wedge (a \vee (b \vee c)) = a & (Abs_\wedge) \\
    a \wedge R &= a \wedge ((a \vee b) \vee c) \\
               &= (a \wedge (a \vee b)) \vee (a \wedge c) & (Dist_\wedge) \\
               &= a \vee (a \wedge c) = a & (Abs_\wedge, Abs_\vee)
\end{align*}
Por tanto, $a \wedge L = a \wedge R$.

2) Comprobamos $\neg a \wedge L = \neg a \wedge R$:
\begin{align*}
    \neg a \wedge L &= \neg a \wedge (a \vee (b \vee c)) \\
                    &= (\neg a \wedge a) \vee (\neg a \wedge (b \vee c)) & (Dist_\wedge) \\
                    &= \bot \vee (\neg a \wedge (b \vee c)) & (Comp_\wedge) \\
                    &= \neg a \wedge (b \vee c) & (ElemNeu_\vee)
\end{align*}
\begin{align*}
    \neg a \wedge R &= \neg a \wedge ((a \vee b) \vee c) \\
                    &= (\neg a \wedge (a \vee b)) \vee (\neg a \wedge c) & (Dist_\wedge) \\
                    &= ((\neg a \wedge a) \vee (\neg a \wedge b)) \vee (\neg a \wedge c) & (Dist_\wedge) \\
                    &= (\bot \vee (\neg a \wedge b)) \vee (\neg a \wedge c) & (Comp_\wedge) \\
                    &= (\neg a \wedge b) \vee (\neg a \wedge c) & (ElemNeu_\vee) \\
                    &= \neg a \wedge (b \vee c) & (Dist_\wedge)
\end{align*}
Como $\neg a \wedge L = \neg a \wedge R$, aplicando el Lema concluimos que $L = R$, es decir, $a \vee (b \vee c) = (a \vee b) \vee c$.
\end{proof}

\begin{proof}[Demostración de $a \wedge (b \wedge c) = (a \wedge b) \wedge c$ (Dual mediante De Morgan)]
Al haber demostrado previamente las leyes de De Morgan, podemos probar la asociatividad del ínfimo de forma directa y elegante sin necesidad de repetir la manipulación algebraica de la demostración dual:
\begin{align*}
    a \wedge (b \wedge c) &= \neg (\neg a \vee \neg (b \wedge c)) & (Mor_\wedge \text{ y } Comp_{inv}) \\
                          &= \neg (\neg a \vee (\neg b \vee \neg c)) & (Mor_\wedge) \\
                          &= \neg ((\neg a \vee \neg b) \vee \neg c) & (Asoc_\vee \text{ demostrada arriba}) \\
                          &= \neg (\neg (a \wedge b) \vee \neg c) & (Mor_\wedge) \\
                          &= (a \wedge b) \wedge c & (Mor_\wedge \text{ y } Comp_{inv})
\end{align*}
\end{proof}
}

\section{Generalización a $n$ variables}

Habiendo demostrado la asociatividad ($Asoc_\vee$ y $Asoc_\wedge$) de las operaciones fundamentales del álgebra de Boole, el orden en el que se agrupan las variables al aplicar consecutivamente una misma operación resulta irrelevante. Esto nos permite prescindir de los paréntesis y extender de forma natural las operaciones binarias a un número arbitrario $n$ de operandos.

\begin{definicion}{Disyunción (Supremo) de $n$ variables}{or_n_variables}
La disyunción múltiple de $n$ variables, denotada de forma compacta mediante el operador $\bigvee$, se define como la aplicación sucesiva de la operación $\vee$:
\[ \bigvee_{i=1}^n x_i \triangleq x_1 \vee x_2 \vee \dots \vee x_n \]
\end{definicion}

\begin{definicion}{Conjunción (Ínfimo) de $n$ variables}{and_n_variables}
De manera análoga, la conjunción múltiple de $n$ variables, denotada mediante el operador $\bigwedge$, se define como la aplicación sucesiva de la operación $\wedge$:
\[ \bigwedge_{i=1}^n x_i \triangleq x_1 \wedge x_2 \wedge \dots \wedge x_n \]
\end{definicion}

La existencia de estas operaciones múltiples bien definidas es un pilar fundamental para desarrollar formas canónicas (como la suma de productos o producto de sumas) y, como veremos a continuación, servirá de base para extender el número de entradas de los operadores derivados.


\section{Operadores Derivados: NAND, NOR, XOR y XNOR}
Las ecuaciones obtenidas a partir de las Leyes de De Morgan demuestran que las operaciones básicas $\vee$ y $\wedge$ pueden ser expresadas íntegramente en términos de la negación de su operación dual. Esto motiva la definición de varios operadores lógicos fundamentales en sistemas digitales. Dos de ellos (NAND y NOR) son de gran relevancia por ser funcionalmente completos por sí solos, mientras que otros dos (XOR y XNOR) son esenciales para funciones aritméticas y de comprobación de paridad:

\begin{definicion}{Operador NAND (Barra de Sheffer)}{nand}
Denotado clásicamente con una flecha hacia arriba ($\uparrow$), se define como la negación del ínfimo.
\[ a \uparrow b \triangleq \neg (a \wedge b) = \neg a \vee \neg b \]
\end{definicion}

\begin{definicion}{Operador NOR (Flecha de Peirce)}{nor}
Denotado con una flecha hacia abajo ($\downarrow$), se define como la negación del supremo.
\[ a \downarrow b \triangleq \neg (a \vee b) = \neg a \wedge \neg b \]
\end{definicion}

\begin{definicion}{Operador XOR (O-exclusiva)}{xor}
Denotado con el símbolo de suma exclusiva ($\oplus$), evalúa a $\top$ cuando exactamente uno de los operandos es $\top$ y el otro $\bot$.
\[ a \oplus b \triangleq (a \wedge \neg b) \vee (\neg a \wedge b) \]
\end{definicion}

\begin{definicion}{Operador XNOR (No-O-exclusiva o Equivalencia)}{xnor}
Denotado frecuentemente con $\odot$ o $\leftrightarrow$, es la negación de la operación XOR y evalúa a $\top$ cuando ambos operandos son idénticos.
\[ a \odot b \triangleq \neg (a \oplus b) = (a \wedge b) \vee (\neg a \wedge \neg b) \]
\end{definicion}

\begin{definicion}{Generalización a $n$ variables de NAND y NOR}{gen_nand_nor}
A diferencia de los operadores $\vee$, $\wedge$ y $\oplus$, los operadores NAND ($\uparrow$) y NOR ($\downarrow$) \textbf{no son asociativos}. Sin embargo, debido a su inmensa importancia práctica en la construcción de circuitos digitales, se define convencionalmente su generalización a $n$ variables como la negación de la conjunción o disyunción múltiple, respectivamente:
\[ \text{NAND}(x_1, x_2, \dots, x_n) \triangleq \neg \left( \bigwedge_{i=1}^n x_i \right) \]
\[ \text{NOR}(x_1, x_2, \dots, x_n) \triangleq \neg \left( \bigvee_{i=1}^n x_i \right) \]
\end{definicion}

\section{Comportamiento de los Operadores Derivados}
Los operadores NAND ($\uparrow$) y NOR ($\downarrow$) presentan una serie de propiedades algebraicas particulares. Al ser operadores funcionalmente completos, permiten expresar cualquier otra operación booleana utilizando exclusivamente uno de ellos.

\begin{teorema}{Idempotencia cruzada (Generación de NOT)}{idemp_cruzada}
Operar un elemento consigo mismo usando NAND o NOR equivale a su complemento (negación):
\begin{align*}
    \forall a \in \mathbb{B}, \quad a \uparrow a &= \neg a \\
    \forall a \in \mathbb{B}, \quad a \downarrow a &= \neg a
\end{align*}
\end{teorema}
\desplegable{proof_idemp_cruzada}{
\begin{proof}[Demostración]
Para la operación NAND:
\begin{align*}
    a \uparrow a &\triangleq \neg (a \wedge a) & (\text{Definicion de NAND}) \\
                 &= \neg a & (Idemp_\wedge)
\end{align*}
Para la operación NOR:
\begin{align*}
    a \downarrow a &\triangleq \neg (a \vee a) & (\text{Definicion de NOR}) \\
                   &= \neg a & (Idemp_\vee)
\end{align*}
\end{proof}
}

\begin{teorema}{Generación del Ínfimo y Supremo (AND y OR)}{gen_inf_sup}
A partir de la propiedad anterior, podemos recuperar las operaciones básicas anidando las puertas consigo mismas:
\begin{align*}
    \forall a, b \in \mathbb{B}, \quad a \wedge b &= \neg (a \uparrow b) = (a \uparrow b) \uparrow (a \uparrow b) \\
    \forall a, b \in \mathbb{B}, \quad a \vee b &= \neg (a \downarrow b) = (a \downarrow b) \downarrow (a \downarrow b)
\end{align*}
\end{teorema}
\desplegable{proof_gen_inf_sup}{
\begin{proof}[Demostración]
Para la generación del ínfimo (AND):
\begin{align*}
    \neg(a \uparrow b) &\triangleq \neg(\neg(a \wedge b)) & (\text{Definicion de NAND}) \\
                       &= a \wedge b & (Involucion)
\end{align*}
Además, por la idempotencia cruzada demostrada anteriormente, $x \uparrow x = \neg x$, por tanto:
\[ \neg(a \uparrow b) = (a \uparrow b) \uparrow (a \uparrow b) \]
La demostración para la generación del supremo (OR) es idéntica por dualidad.
\end{proof}
}

\begin{teorema}{Generación cruzada (Leyes de De Morgan para NAND/NOR)}{gen_cruzada}
Podemos generar la operación opuesta (supremo desde NAND, e ínfimo desde NOR) negando previamente las entradas:
\begin{align*}
    \forall a, b \in \mathbb{B}, \quad a \vee b &= (\neg a) \uparrow (\neg b) = (a \uparrow a) \uparrow (b \uparrow b) \\
    \forall a, b \in \mathbb{B}, \quad a \wedge b &= (\neg a) \downarrow (\neg b) = (a \downarrow a) \downarrow (b \downarrow b)
\end{align*}
\end{teorema}
\desplegable{proof_gen_cruzada}{
\begin{proof}[Demostración]
Para el supremo (OR):
\begin{align*}
    (\neg a) \uparrow (\neg b) &\triangleq \neg (\neg a \wedge \neg b) & (\text{Definicion de NAND}) \\
                               &= \neg (\neg (a \vee b)) & (Mor_\vee) \\
                               &= a \vee b & (Involucion)
\end{align*}
La demostración para el ínfimo (AND) sigue los mismos pasos de manera dual, aplicando $Mor_\wedge$.
\end{proof}
}

\begin{teorema}{Conmutatividad}{conmut_deriv}
Al igual que sus operaciones base, ambos operadores son perfectamente conmutativos:
\begin{align*}
    \forall a, b \in \mathbb{B}, \quad a \uparrow b &= b \uparrow a \\
    \forall a, b \in \mathbb{B}, \quad a \downarrow b &= b \downarrow a
\end{align*}
\end{teorema}
\desplegable{proof_conmut_deriv}{
\begin{proof}[Demostración]
Para la operación NAND:
\begin{align*}
    a \uparrow b &\triangleq \neg (a \wedge b) & (\text{Definicion de NAND}) \\
                 &= \neg (b \wedge a) & (Comm_\wedge) \\
                 &\triangleq b \uparrow a & (\text{Definicion de NAND})
\end{align*}
Para la operación NOR, es análogo aplicando $Comm_\vee$.
\end{proof}
}

\begin{teorema}{Inexistencia de Elemento Neutro}{no_neutro}
No existe ningún elemento neutro para las operaciones NAND ni NOR en un álgebra de Boole general.
\end{teorema}
\desplegable{proof_no_neutro}{
\begin{proof}[Demostración]
Si existiera un neutro $e$ para la operación NAND, debería cumplirse que $\forall a, a \uparrow e = a$, es decir, $\neg(a \wedge e) = a$. Si probamos con $e=\top$, obtenemos $\neg a = a$, lo cual obliga al colapso en un álgebra trivial. Si probamos con $e=\bot$, obtenemos $\neg \bot = a \implies \top = a$, lo cual obviamente no se cumple para cualquier elemento $a$. Lo mismo aplica a la operación NOR.
\end{proof}
}

\begin{teorema}{Comportamiento con las constantes (Fijación y Absorción)}{constantes_deriv}
Fijar una constante específica en uno de los operandos genera directamente la negación, mientras que usar la constante opuesta actúa como un pseudo-elemento absorbente (devolviendo un valor constante inalterable por $a$):
\begin{align*}
    \text{Inversión: } & a \uparrow \top = \neg a \qquad & a \downarrow \bot &= \neg a \\
    \text{Absorción: } & a \uparrow \bot = \top \qquad & a \downarrow \top &= \bot
\end{align*}
\end{teorema}
\desplegable{proof_constantes_deriv}{
\begin{proof}[Demostración]
Para la inversión:
\begin{align*}
    a \uparrow \top &= \neg(a \wedge \top) = \neg a & (ElemNeu_\wedge) \\
    a \downarrow \bot &= \neg(a \vee \bot) = \neg a & (ElemNeu_\vee)
\end{align*}
Para la absorción:
\begin{align*}
    a \uparrow \bot &= \neg(a \wedge \bot) = \neg \bot = \top & (Abs_\bot) \\
    a \downarrow \top &= \neg(a \vee \top) = \neg \top = \bot & (Abs_\top)
\end{align*}
\end{proof}
}

\begin{teorema}{Ausencia de Asociatividad}{no_asoc_deriv}
A diferencia del supremo ($\vee$) y el ínfimo ($\wedge$), las operaciones NAND y NOR son positivamente NO asociativas:
\begin{align*}
    (a \uparrow b) \uparrow c &\neq a \uparrow (b \uparrow c) \\
    (a \downarrow b) \downarrow c &\neq a \downarrow (b \downarrow c)
\end{align*}
\end{teorema}
\desplegable{proof_no_asoc_deriv}{
\begin{proof}[Demostración]
Desarrollando el lado izquierdo para NAND:
\[ (a \uparrow b) \uparrow c = \neg ( (\neg (a \wedge b)) \wedge c ) = (a \wedge b) \vee \neg c \]
Desarrollando el lado derecho:
\[ a \uparrow (b \uparrow c) = \neg ( a \wedge (\neg (b \wedge c)) ) = \neg a \vee (b \wedge c) \]
Resulta evidente que $(a \wedge b) \vee \neg c \neq \neg a \vee (b \wedge c)$ para combinaciones arbitrarias de variables. El mismo razonamiento aplica de manera estricta y análoga para demostrar la carencia de asociatividad en la operación NOR ($\downarrow$).
\end{proof}
}

\begin{definicion}{NAND y NOR de $n$ entradas}{def_n_entradas}
Dado que las puertas NAND y NOR físicas a menudo tienen más de dos entradas, se definen algebraicamente para múltiples entradas como la negación de la conjunción o disyunción de todas ellas:
\begin{align*}
    \uparrow(x_1, x_2, \dots, x_n) &\triangleq \neg \left( \bigwedge_{i=1}^n x_i \right) \\
    \downarrow(x_1, x_2, \dots, x_n) &\triangleq \neg \left( \bigvee_{i=1}^n x_i \right)
\end{align*}
En particular, para el caso de 3 entradas que estudiaremos a continuación:
\begin{align*}
    \uparrow(a,b,c) &\triangleq \neg(a \wedge b \wedge c) \\
    \downarrow(a,b,c) &\triangleq \neg(a \vee b \vee c)
\end{align*}
\end{definicion}

\begin{teorema}{NAND/NOR múltiple vs agrupación binaria}{multiple_vs_binaria}
Como consecuencia directa de su falta de asociatividad, una operación NAND o NOR de 3 entradas no es equivalente a la agrupación secuencial en cascada de operaciones de 2 entradas:
\begin{align*}
    \uparrow(a,b,c) &\neq (a \uparrow b) \uparrow c \qquad \text{y} \qquad \uparrow(a,b,c) \neq a \uparrow (b \uparrow c) \\
    \downarrow(a,b,c) &\neq (a \downarrow b) \downarrow c \qquad \text{y} \qquad \downarrow(a,b,c) \neq a \downarrow (b \downarrow c)
\end{align*}
\end{teorema}
\desplegable{proof_multiple_vs_binaria}{
\begin{proof}[Demostración]
Para la NAND de 3 entradas tenemos, por definición y leyes de De Morgan:
\[ \uparrow(a,b,c) \triangleq \neg(a \wedge b \wedge c) = \neg a \vee \neg b \vee \neg c \]
Sin embargo, la agrupación de dos en dos evaluada anteriormente daba:
\[ (a \uparrow b) \uparrow c = (a \wedge b) \vee \neg c \]
Evidentemente, $\neg a \vee \neg b \vee \neg c \neq (a \wedge b) \vee \neg c$. Lo mismo aplica a las agrupaciones derechas y a las operaciones NOR equivalentes.
\end{proof}
}

\begin{teorema}{Extensión del Principio de Dualidad (NAND y NOR)}{dualidad_nand_nor}
La inclusión de los operadores derivados expande el Principio de Dualidad establecido en los postulados iniciales. La expresión dual de cualquier teorema o identidad que contenga operaciones NAND o NOR se obtiene intercambiando los operadores $\uparrow$ y $\downarrow$ (además de los ya conocidos $\vee \leftrightarrow \wedge$ y $\bot \leftrightarrow \top$).
\end{teorema}
\subsection{Comportamiento de los Operadores XOR y XNOR}
A diferencia de los operadores NAND y NOR, que destacan por su universalidad funcional pero carecen de propiedades algebraicas deseables (como asociatividad o elemento neutro), los operadores XOR ($\oplus$) y XNOR ($\odot$) exhiben una rica estructura algebraica. A continuación demostraremos estas propiedades.

\begin{teorema}{Conmutatividad}{conmut_xor}
Ambos operadores son conmutativos:
\[ a \oplus b = b \oplus a \qquad \text{y} \qquad a \odot b = b \odot a \]
\end{teorema}
\desplegable{proof_conmut_xor}{
\begin{proof}[Demostración]
Para XOR:
\begin{align*}
    a \oplus b &\triangleq (a \wedge \neg b) \vee (\neg a \wedge b) & (\text{Definición de XOR}) \\
               &= (\neg a \wedge b) \vee (a \wedge \neg b) & (Comm_\vee) \\
               &= (b \wedge \neg a) \vee (\neg b \wedge a) & (Comm_\wedge) \\
               &\triangleq b \oplus a & (\text{Definición de XOR})
\end{align*}
La demostración para XNOR sigue pasos idénticos aprovechando la conmutatividad del ínfimo y el supremo.
\end{proof}
}

\begin{teorema}{Elementos Neutros e Inversores}{neutro_xor}
El elemento $\bot$ actúa como neutro para la XOR, y $\top$ actúa como inversor. De manera dual, $\top$ es el neutro de la XNOR, y $\bot$ actúa como inversor:
\begin{align*}
    a \oplus \bot &= a & a \oplus \top &= \neg a \\
    a \odot \top &= a & a \odot \bot &= \neg a
\end{align*}
\end{teorema}
\desplegable{proof_neutro_xor}{
\begin{proof}[Demostración para XOR]
Para el elemento $\bot$ (neutro):
\begin{align*}
    a \oplus \bot &\triangleq (a \wedge \neg \bot) \vee (\neg a \wedge \bot) & (\text{Definición}) \\
                  &= (a \wedge \top) \vee \bot & (Comp_{inv} \text{ y } Fij_\wedge) \\
                  &= a \vee \bot & (ElemNeu_\wedge) \\
                  &= a & (ElemNeu_\vee)
\end{align*}
Para el elemento $\top$ (inversor):
\begin{align*}
    a \oplus \top &\triangleq (a \wedge \neg \top) \vee (\neg a \wedge \top) & (\text{Definición}) \\
                  &= (a \wedge \bot) \vee \neg a & (Comp_{inv} \text{ y } ElemNeu_\wedge) \\
                  &= \bot \vee \neg a & (Fij_\wedge) \\
                  &= \neg a & (ElemNeu_\vee)
\end{align*}
Las pruebas para XNOR son totalmente duales.
\end{proof}
}

\begin{teorema}{Elemento Inverso de sí mismo (Grupo Abeliano)}{idemp_nula_xor}
La combinación de un elemento consigo mismo produce una anulación (devuelve el neutro de la operación correspondiente), actuando cada elemento como su propio inverso:
\[ a \oplus a = \bot \qquad \text{y} \qquad a \odot a = \top \]
\end{teorema}
\desplegable{proof_idemp_nula_xor}{
\begin{proof}[Demostración]
Para XOR:
\begin{align*}
    a \oplus a &\triangleq (a \wedge \neg a) \vee (\neg a \wedge a) \\
               &= \bot \vee \bot & (Comp_\wedge \text{ y } Comm_\wedge) \\
               &= \bot & (Idemp_\vee)
\end{align*}
Para XNOR:
\begin{align*}
    a \odot a &\triangleq (a \wedge a) \vee (\neg a \wedge \neg a) \\
              &= a \vee \neg a & (Idemp_\wedge) \\
              &= \top & (Comp_\vee)
\end{align*}
\end{proof}
}

\begin{teorema}{Propiedades de Negación}{negacion_xor}
Negar cualquiera de las entradas de forma independiente equivale a negar la operación completa, lo que a su vez alterna entre XOR y XNOR:
\begin{align*}
    \neg (a \oplus b) &= \neg a \oplus b = a \oplus \neg b = a \odot b \\
    \neg (a \odot b) &= \neg a \odot b = a \odot \neg b = a \oplus b
\end{align*}
\end{teorema}
\desplegable{proof_negacion_xor}{
\begin{proof}[Demostración de $a \oplus \neg b = \neg(a \oplus b)$]
\begin{align*}
    a \oplus \neg b &\triangleq (a \wedge \neg(\neg b)) \vee (\neg a \wedge \neg b) \\
                    &= (a \wedge b) \vee (\neg a \wedge \neg b) & (Involucion) \\
                    &\triangleq a \odot b & (\text{Definición de XNOR})
\end{align*}
Sabiendo por definición que $a \odot b \triangleq \neg(a \oplus b)$, se concluye de forma inmediata que $a \oplus \neg b = \neg (a \oplus b)$.
\end{proof}
}

\begin{teorema}{Asociatividad}{asoc_xor}
Tanto XOR como XNOR son operadores algebraicamente asociativos:
\[ (a \oplus b) \oplus c = a \oplus (b \oplus c) \qquad \text{y} \qquad (a \odot b) \odot c = a \odot (b \odot c) \]
\end{teorema}
\desplegable{proof_asoc_xor}{
\begin{proof}[Demostración para XOR]
Primero evaluaremos el miembro izquierdo $(a \oplus b) \oplus c$. Llamaremos $X = a \oplus b$.
\begin{align*}
    X \oplus c &\triangleq (X \wedge \neg c) \vee (\neg X \wedge c) \\
               &= \big( ((a \wedge \neg b) \vee (\neg a \wedge b)) \wedge \neg c \big) \vee \big( \neg ((a \wedge \neg b) \vee (\neg a \wedge b)) \wedge c \big) \\
               &= \big( ((a \wedge \neg b) \vee (\neg a \wedge b)) \wedge \neg c \big) \vee \big( (a \odot b) \wedge c \big) \quad (\text{Definición de XNOR}) \\
               &= \big( (a \wedge \neg b \wedge \neg c) \vee (\neg a \wedge b \wedge \neg c) \big) \vee \big( ((a \wedge b) \vee (\neg a \wedge \neg b)) \wedge c \big) \quad (Dist_\wedge) \\
               &= (a \wedge \neg b \wedge \neg c) \vee (\neg a \wedge b \wedge \neg c) \\
               &\quad \vee (a \wedge b \wedge c) \vee (\neg a \wedge \neg b \wedge c) \quad (Dist_\wedge)
\end{align*}
Ahora evaluaremos el miembro derecho $a \oplus (b \oplus c)$. Llamaremos $Y = b \oplus c$.
\begin{align*}
    a \oplus Y &\triangleq (a \wedge \neg Y) \vee (\neg a \wedge Y) \\
               &= (a \wedge (b \odot c)) \vee (\neg a \wedge ((b \wedge \neg c) \vee (\neg b \wedge c))) \\
               &= (a \wedge ((b \wedge c) \vee (\neg b \wedge \neg c))) \vee (\neg a \wedge b \wedge \neg c) \vee (\neg a \wedge \neg b \wedge c) \\
               &= (a \wedge b \wedge c) \vee (a \wedge \neg b \wedge \neg c) \\
               &\quad \vee (\neg a \wedge b \wedge \neg c) \vee (\neg a \wedge \neg b \wedge c)
\end{align*}
Como podemos observar, ambas expansiones resultan exactamente en los mismos cuatro minitérminos. Reordenándolos por conmutatividad ($Comm_\vee$) demostramos que son idénticos.
\end{proof}
}

\begin{teorema}{Generalización n-aria (XOR y XNOR)}{gen_xor}
Dado que ambos operadores han demostrado ser asociativos y conmutativos, es posible omitir los paréntesis y generalizar la operación a un número arbitrario $n$ de variables. Se denotan mediante los operadores de sumatoria y productorio modificados:
\[ \bigoplus_{i=1}^{n} x_i = x_1 \oplus x_2 \oplus \dots \oplus x_n \]
\[ \bigodot_{i=1}^{n} x_i = x_1 \odot x_2 \odot \dots \odot x_n \]
\end{teorema}

\begin{teorema}{Distributividad (AND sobre XOR y OR sobre XNOR)}{dist_xor}
El producto (AND) se distribuye sobre la suma exclusiva (XOR), y dualmente, la suma (OR) se distribuye sobre la equivalencia (XNOR):
\[ a \wedge (b \oplus c) = (a \wedge b) \oplus (a \wedge c) \]
\[ a \vee (b \odot c) = (a \vee b) \odot (a \vee c) \]
\end{teorema}
\desplegable{proof_dist_xor}{
\begin{proof}[Demostración de AND sobre XOR]
Desarrollando el lado derecho (RHS):
\begin{align*}
    (a \wedge b) \oplus (a \wedge c) &\triangleq ((a \wedge b) \wedge \neg(a \wedge c)) \vee (\neg(a \wedge b) \wedge (a \wedge c)) \\
                                     &= (a \wedge b \wedge (\neg a \vee \neg c)) \vee ((\neg a \vee \neg b) \wedge a \wedge c) \quad (Mor_\wedge) \\
                                     &= ((a \wedge b \wedge \neg a) \vee (a \wedge b \wedge \neg c)) \\
                                     &\quad \vee ((a \wedge c \wedge \neg a) \vee (a \wedge c \wedge \neg b)) \quad (Dist_\wedge) \\
                                     &= (\bot \vee (a \wedge b \wedge \neg c)) \vee (\bot \vee (a \wedge \neg b \wedge c)) \quad (Comp_\wedge \text{ y } Fij_\wedge) \\
                                     &= (a \wedge b \wedge \neg c) \vee (a \wedge \neg b \wedge c) \quad (ElemNeu_\vee) \\
                                     &= a \wedge ((b \wedge \neg c) \vee (\neg b \wedge c)) \quad (Dist_\wedge \text{ a la inversa}) \\
                                     &\triangleq a \wedge (b \oplus c) \quad (\text{Def. XOR})
\end{align*}
Lo cual demuestra la igualdad. La demostración de la distributividad para XNOR es rigurosamente dual.
\end{proof}
}

\begin{teorema}{Extensión del Principio de Dualidad (XOR y XNOR)}{dualidad_xor_xnor}
De forma análoga a la relación entre NAND y NOR, los operadores XOR y XNOR son mutuamente duales. Para obtener la expresión dual de cualquier proposición que involucre estos operadores, se deben intercambiar $\oplus$ y $\odot$, manteniendo las reglas de dualidad estándar para los demás elementos y constantes.
\end{teorema}

\section{Instanciación en Álgebras Finitas: Trivial y Bivaluada}

La teoría desarrollada en los capítulos anteriores es aplicable a cualquier álgebra de Boole, sin importar el número de elementos que contenga el conjunto $B$ (siempre que se cumplan los postulados de Huntington). Sin embargo, existen dos álgebras finitas de interés particular por su extrema simplicidad y su aplicación directa en la teoría de circuitos.

\subsection{Álgebra Trivial ($|B| = 1$)}
Si definimos el conjunto soporte con un único elemento, $B = \{ c \}$, nos encontramos ante el álgebra de Boole trivial o degenerada. 

Dado que los postulados de Huntington (Postulado 2) exigen la existencia de un elemento neutro para la disyunción ($\bot \in B$) y otro para la conjunción ($\top \in B$), y puesto que el conjunto solo contiene un único elemento, estos deben forzosamente coincidir:
\[ \bot = \top = c \]

Al instanciar cualquier operación definida sobre este conjunto, los resultados siempre evalúan a dicha constante $c$. 
\begin{itemize}
    \item \textbf{Negación:} Por el postulado del complemento, $c \vee \neg c = c$ y $c \wedge \neg c = c$, lo que implica que $\neg c = c$.
    \item \textbf{Disyunción y Conjunción:} Por la propiedad de idempotencia, $c \vee c = c$ y $c \wedge c = c$.
    \item \textbf{Operadores Derivados:} Por definición, $c \uparrow c = \neg (c \wedge c) = \neg c = c$. Lo mismo sucede con el resto de operadores.
\end{itemize}

Visualizar esto en tablas de operación (comúnmente conocidas como tablas de verdad) resulta en estructuras degeneradas de una sola celda, donde $\circ \in \{ \vee, \wedge, \uparrow, \downarrow, \oplus, \odot \}$:

\begin{center}
\begin{tabular}{c|c}
$a$ & $\neg a$ \\
\hline
$c$ & $c$
\end{tabular}
\qquad
\begin{tabular}{cc|c}
$a$ & $b$ & $a \circ b$ \\
\hline
$c$ & $c$ & $c$
\end{tabular}
\end{center}

\subsection{Álgebra Bivaluada ($|B| = 2$)}
El caso más importante para la ingeniería es el álgebra de Boole bivaluada, donde el conjunto soporte consta exactamente de los dos elementos garantizados por los postulados: el neutro disyuntivo y el neutro conjuntivo. 
\[ B = \{ \bot, \top \} \]
(Asumiendo lógicamente que $\bot \neq \top$).

Procederemos a deducir el comportamiento (las tablas de operación) instanciando los teoremas y postulados axiomáticos en estos dos únicos valores.

\subsubsection{Operaciones Básicas ($\neg$, $\vee$, $\wedge$)}

\textbf{1. Negación (Operación unaria $\neg$)} \\
El Postulado 5 (Complemento) exige que:
\[ \bot \vee \neg \bot = \top \quad \text{y} \quad \top \vee \neg \top = \top \]
Al existir solo dos elementos en el conjunto, el único valor que sumado a $\bot$ (que es el neutro disyuntivo, por lo que no altera el resultado) da $\top$, es el propio $\top$. Por lo tanto, deducimos que $\neg \bot = \top$. De igual manera, por dualidad, $\neg \top = \bot$.

\textbf{2. Disyunción (Operación binaria $\vee$)} \\
Calculamos los cuatro casos posibles instanciando los valores:
\begin{itemize}
    \item $\bot \vee \bot = \bot$ (Por Idempotencia, Teorema 1).
    \item $\bot \vee \top = \top$ (Por ser $\bot$ el elemento neutro, Postulado 2a).
    \item $\top \vee \bot = \top$ (Por Conmutatividad, Postulado 3a).
    \item $\top \vee \top = \top$ (Por Idempotencia, Teorema 1).
\end{itemize}

\textbf{3. Conjunción (Operación binaria $\wedge$)} \\
Análogamente:
\begin{itemize}
    \item $\top \wedge \top = \top$ (Por Idempotencia, Teorema 1).
    \item $\top \wedge \bot = \bot$ (Por ser $\top$ el elemento neutro, Postulado 2b).
    \item $\bot \wedge \top = \bot$ (Por Conmutatividad, Postulado 3b).
    \item $\bot \wedge \bot = \bot$ (Por Idempotencia, Teorema 1).
\end{itemize}

\subsubsection{Operadores Derivados ($\uparrow, \downarrow, \oplus, \odot$)}
Podemos obtener las tablas de los operadores derivados aplicando directamente sus definiciones algebraicas sobre las tablas básicas ya obtenidas:

\textbf{1. NAND y NOR} \\
Dado que $a \uparrow b = \neg(a \wedge b)$ y $a \downarrow b = \neg(a \vee b)$, los resultados consisten simplemente en aplicar el operador complemento ($\neg$) a las tablas de conjunción y disyunción calculadas previamente.

\textbf{2. XOR y XNOR} \\
Recordando la definición algebraica $a \oplus b = (a \wedge \neg b) \vee (\neg a \wedge b)$, se puede evaluar caso por caso (por ejemplo, $\top \oplus \bot = (\top \wedge \top) \vee (\bot \wedge \bot) = \top \vee \bot = \top$), pero también podemos usar directamente los teoremas derivados anteriormente:
\begin{itemize}
    \item $a \oplus \bot = a$ (Elemento neutro). Por lo tanto: $\bot \oplus \bot = \bot$, y $\top \oplus \bot = \top$.
    \item $a \oplus \top = \neg a$ (Inversor). Por lo tanto: $\bot \oplus \top = \top$, y $\top \oplus \top = \bot$.
\end{itemize}
Por dualidad, y sabiendo que el XNOR es la negación del XOR, se obtiene trivialmente que $a \odot b = \neg(a \oplus b)$.

\subsubsection{Resumen: Tablas de Operación Bivaluadas}
A continuación, presentamos la consolidación matricial de todas las operaciones deducidas, conformando las tablas de verdad definitivas del álgebra de dos valores:

\begin{table}[h!]
\centering
\renewcommand{\arraystretch}{1.2}
\begin{tabular}{c||c}
\hline
$a$ & $\neg a$ \\
\hline\hline
$\bot$ & $\top$ \\
$\top$ & $\bot$ \\
\hline
\end{tabular}
\qquad
\begin{tabular}{cc||c|c|c|c|c|c}
\hline
$a$ & $b$ & $a \vee b$ & $a \wedge b$ & $a \uparrow b$ & $a \downarrow b$ & $a \oplus b$ & $a \odot b$ \\
\hline\hline
$\bot$ & $\bot$ & $\bot$ & $\bot$ & $\top$ & $\top$ & $\bot$ & $\top$ \\
$\bot$ & $\top$ & $\top$ & $\bot$ & $\top$ & $\bot$ & $\top$ & $\bot$ \\
$\top$ & $\bot$ & $\top$ & $\bot$ & $\top$ & $\bot$ & $\top$ & $\bot$ \\
$\top$ & $\top$ & $\top$ & $\top$ & $\bot$ & $\bot$ & $\bot$ & $\top$ \\
\hline
\end{tabular}
\caption{Tablas de Verdad consolidadas para las Operaciones del Álgebra de Boole de 2 elementos.}
\label{tab:tablas_bivaluadas}
\end{table}

\section{Conclusiones y Transición a la Lógica Digital}

Tras haber establecido formalmente la estructura matemática del álgebra de Boole a partir de los postulados de Huntington, y haber demostrado rigurosamente sus propiedades fundamentales operando con la signatura clásica de la teoría de retículos ($\vee, \wedge, \bot, \top$), estamos en disposición de dar el salto al dominio de la ingeniería.

En la electrónica digital, el interés recae de forma exclusiva sobre un modelo concreto de álgebra de Boole: el \textbf{Álgebra de Conmutación de Shannon}. Esta es el álgebra de Boole más sencilla posible, cuyo conjunto subyacente consta únicamente de dos elementos, $\mathbb{B}_2 = \{0, 1\}$. 

A pesar de su aparente simplicidad, el álgebra de $\mathbb{B}_2$ cumple estrictamente todos los postulados de Huntington (y por ende, todos los teoremas que hemos derivado) y está contenida formalmente como subestructura en cualquier otra álgebra de Boole más compleja.

Para adecuar nuestra matemática al diseño de circuitos y sistemas digitales, adoptaremos a partir de ahora la \textbf{notación ingenieril}, realizando el siguiente isomorfismo simbólico sobre nuestras operaciones y constantes:
\begin{itemize}
    \item \textbf{Constantes lógicas:} El elemento mínimo $\bot$ (falso) se denotará como \textbf{0} (ó nivel bajo de tensión, $L$). El elemento máximo $\top$ (verdadero) se denotará como \textbf{1} (ó nivel alto de tensión, $H$).
    \item \textbf{Supremo (Join / Disyunción):} La operación $\vee$ se denotará mediante el operador suma $\mathbf{+}$. En circuitos lógicos, implementa la puerta \textbf{OR}.
    \item \textbf{Ínfimo (Meet / Conjunción):} La operación $\wedge$ se denotará mediante el operador producto $\mathbf{\cdot}$ (frecuentemente omitido, escribiendo $ab$ en lugar de $a \cdot b$). Implementa la puerta \textbf{AND}.
    \item \textbf{Complemento (Negación):} La operación de complemento $\neg a$ se denotará convencionalmente colocando una barra superior sobre la variable, $\mathbf{\overline{a}}$, o mediante una comilla $\mathbf{a'}$. Implementa la puerta \textbf{NOT} (inversor).
    \item \textbf{Operadores Derivados (NAND y NOR):} Las operaciones $\uparrow$ y $\downarrow$ mantienen sus símbolos, o bien se expresan directamente como el complemento del producto o de la suma ($\overline{a \cdot b}$, $\overline{a+b}$). Representan las puertas universales \textbf{NAND} y \textbf{NOR}, fundamentales en el diseño de circuitos integrados.
    \item \textbf{Suma Exclusiva y Equivalencia (XOR y XNOR):} Las operaciones introducidas como suma exclusiva y equivalencia lógica se denotan mediante $\mathbf{\oplus}$ y $\mathbf{\odot}$. Representan las puertas \textbf{XOR} (útiles en sumadores o detectores de paridad) y \textbf{XNOR} (comparadores de igualdad).
    \item \textbf{Operadores n-arios:} Las versiones generalizadas para múltiples variables formarán estructuras de puertas lógicas de $n$ entradas. Se denotarán mediante los operadores $\sum$ (puerta OR de $n$ entradas), $\prod$ (puerta AND de $n$ entradas), $\bigoplus$ (puerta XOR de $n$ entradas) y $\bigodot$ (puerta XNOR de $n$ entradas).
\end{itemize}

Esta notación algebraica clásica resulta mucho más ágil y familiar para la manipulación y simplificación de funciones lógicas complejas. Así, teoremas como el de la distributividad se reescriben de forma natural como $a \cdot (b + c) = (a \cdot b) + (a \cdot c)$, y las leyes de De Morgan cobran su célebre forma visual:
\[ \overline{a + b} = \overline{a} \cdot \overline{b} \qquad \text{y} \qquad \overline{a \cdot b} = \overline{a} + \overline{b} \]
De igual modo, la estructura de las operaciones derivadas queda plasmada directamente en ecuaciones como $a \oplus b = (a \cdot \overline{b}) + (\overline{a} \cdot b)$.

Además, gracias a nuestra previa instanciación en el álgebra bivaluada ($|B|=2$), sabemos que las tablas de operación algebraicas que hemos deducido analíticamente se corresponden de manera idéntica y biunívoca con las \textbf{tablas de verdad} de las puertas lógicas físicas.

Con estos fundamentos matemáticos sólidamente establecidos, la transición hacia el diseño, análisis y simplificación de circuitos digitales queda completamente justificada y carente de ambigüedades.

\newpage
\section{Anexo: Resumen de Postulados y Teoremas (Notación Ingenieril)}
\makeatletter

En esta sección se recopilan los postulados de Huntington y los teoremas principales derivados, transcritos a la notación propia del álgebra de conmutación y la lógica digital ($+$, $\cdot$, $0$, $1$, $\overline{a}$), concebidos como hoja de referencia rápida.

\subsection*{Pre-Axiomas de la Estructura}

\setcounter{tcb@cnt@preaxioma}{0}
\begin{preaxioma}{Estructura de Conjunto}{esconj_eng}
Se requiere que se defina sobre un conjunto (por ejemplo, $\mathbb{B}_2 = \{0, 1\}$).
\end{preaxioma}

\setcounter{tcb@cnt@preaxioma}{1}
\begin{preaxioma}{Constantes Lógicas}{constantes_eng}
Este conjunto contiene dos constantes fundamentales: $0$ (falso) y $1$ (verdadero).
\end{preaxioma}

\setcounter{tcb@cnt@preaxioma}{2}
\begin{preaxioma}{Operaciones Binarias Internas}{opbinint_eng}
Se definen dos operaciones binarias internas, la suma ($+$) y el producto ($\cdot$):
\begin{align*}
+ &: \mathbb{B}_2 \times \mathbb{B}_2 \to \mathbb{B}_2 \\
\cdot &: \mathbb{B}_2 \times \mathbb{B}_2 \to \mathbb{B}_2
\end{align*}
\end{preaxioma}

\setcounter{tcb@cnt@preaxioma}{4}
\begin{preaxioma}{Existencia y Unicidad de Imagen}{exist_unic_eng}
Para cada par de elementos del conjunto, las operaciones $+$ y $\cdot$ siempre producen un resultado que también pertenece al conjunto, y ese resultado es siempre único.
\end{preaxioma}

\subsection*{Postulados de Huntington}

\setcounter{tcb@cnt@postulado}{0}
\begin{postulado}{Elemento neutro}{neutro_eng}
\[ a + 0 = a \qquad \text{y} \qquad a \cdot 1 = a \]
\end{postulado}

\setcounter{tcb@cnt@postulado}{2}
\begin{postulado}{Conmutatividad}{conmut_eng}
\[ a + b = b + a \qquad \text{y} \qquad a \cdot b = b \cdot a \]
\end{postulado}

\setcounter{tcb@cnt@postulado}{4}
\begin{postulado}{Distributividad}{distrib_eng}
\[ a \cdot (b + c) = (a \cdot b) + (a \cdot c) \qquad \text{y} \qquad a + (b \cdot c) = (a + b) \cdot (a + c) \]
\end{postulado}

\setcounter{tcb@cnt@postulado}{6}
\begin{postulado}{Complementario}{comp_eng}
Para cada elemento $a$, existe un complemento $\overline{a}$ tal que:
\[ a + \overline{a} = 1 \qquad \text{y} \qquad a \cdot \overline{a} = 0 \]
\end{postulado}

\subsection*{Teoremas Fundamentales}

\setcounter{tcb@cnt@teorema}{0}
\begin{teorema}{Unicidad de los elementos neutros}{unicidad_neutros_eng}
El elemento neutro para la suma ($0$) y para el producto ($1$) son únicos.
\end{teorema}

\setcounter{tcb@cnt@teorema}{1}
\begin{teorema}{Idempotencia}{idempotencia_eng}
\[ a + a = a \qquad \text{y} \qquad a \cdot a = a \]
\end{teorema}

\setcounter{tcb@cnt@teorema}{2}
\begin{teorema}{Elementos absorbentes}{absorbentes_eng}
\[ a + 1 = 1 \qquad \text{y} \qquad a \cdot 0 = 0 \]
\end{teorema}

\setcounter{tcb@cnt@teorema}{5}
\begin{teorema}{Propiedades de absorción}{absorcion_eng}
\[ a + (a \cdot b) = a \qquad \text{y} \qquad a \cdot (a + b) = a \]
\end{teorema}

\setcounter{tcb@cnt@teorema}{11}
\begin{teorema}{Leyes de De Morgan}{morgan_eng}
\[ \overline{a + b} = \overline{a} \cdot \overline{b} \qquad \text{y} \qquad \overline{a \cdot b} = \overline{a} + \overline{b} \]
\end{teorema}

\setcounter{tcb@cnt@teorema}{10}
\begin{teorema}{Involución (Doble negación)}{involucion_eng}
\[ \overline{\overline{a}} = a \]
\end{teorema}

\setcounter{tcb@cnt@teorema}{12}
\begin{teorema}{Asociatividad}{asociatividad_eng}
\[ a + (b + c) = (a + b) + c \qquad \text{y} \qquad a \cdot (b \cdot c) = (a \cdot b) \cdot c \]
\end{teorema}

\setcounter{tcb@cnt@teorema}{9}
\begin{teorema}{Unicidad del complemento}{unic_comp_eng}
El complemento $\overline{a}$ de un elemento $a$ es único.
\end{teorema}

\setcounter{tcb@cnt@teorema}{6}
\begin{teorema}{Otras propiedades equivalentes}{otras_prop_eng}
\begin{itemize}
    \item \textbf{Orden de retículo:} $a + b = a \iff a \cdot b = b$
    \item \textbf{Equivalencia de operaciones:} $a + b = a \cdot b \implies a = b$
    \item \textbf{Cancelación:} $(a + b = a + c \text{ y } a \cdot b = a \cdot c) \implies b = c$
\end{itemize}
\end{teorema}

\setcounter{tcb@cnt@definicion}{0}
\begin{definicion}{Generalización a $n$ variables}{gen_n_vars_eng}
Las operaciones disyunción y conjunción pueden extenderse a un número $n$ de variables mediante los símbolos sumatorio y productorio:
\[ \sum_{i=1}^{n} x_i = x_1 + x_2 + \dots + x_n \]
\[ \prod_{i=1}^{n} x_i = x_1 \cdot x_2 \cdot \dots \cdot x_n \]
\end{definicion}

\setcounter{tcb@cnt@teorema}{3}
\begin{teorema}{Casos de Álgebra Trivial}{trivial_eng}
Si $0 = 1$, o si existe algún elemento tal que $\overline{a} = a$, entonces el álgebra contiene un único elemento (álgebra trivial).
\end{teorema}

\subsection*{Comportamiento de Operadores Derivados}

\setcounter{tcb@cnt@teorema}{13}
\begin{teorema}{Idempotencia cruzada (NAND/NOR)}{idemp_cruzada_eng}
\[ a \uparrow a = \overline{a} \qquad \text{y} \qquad a \downarrow a = \overline{a} \]
\end{teorema}

\setcounter{tcb@cnt@teorema}{14}
\begin{teorema}{Generación de AND y OR}{gen_inf_sup_eng}
\[ a \cdot b = \overline{a \uparrow b} = (a \uparrow b) \uparrow (a \uparrow b) \qquad \text{y} \qquad a + b = \overline{a \downarrow b} = (a \downarrow b) \downarrow (a \downarrow b) \]
\end{teorema}

\setcounter{tcb@cnt@teorema}{15}
\begin{teorema}{Generación cruzada}{gen_cruzada_eng}
\[ a + b = \overline{a} \uparrow \overline{b} = (a \uparrow a) \uparrow (b \uparrow b) \qquad \text{y} \qquad a \cdot b = \overline{a} \downarrow \overline{b} = (a \downarrow a) \downarrow (b \downarrow b) \]
\end{teorema}

\setcounter{tcb@cnt@teorema}{16}
\begin{teorema}{Conmutatividad}{conmut_deriv_eng}
\[ a \uparrow b = b \uparrow a \qquad \text{y} \qquad a \downarrow b = b \downarrow a \]
\end{teorema}

\setcounter{tcb@cnt@teorema}{18}
\begin{teorema}{Comportamiento con las constantes}{constantes_deriv_eng}
\begin{align*}
    a \uparrow 1 &= \overline{a} \qquad & a \downarrow 0 &= \overline{a} \\
    a \uparrow 0 &= 1 \qquad & a \downarrow 1 &= 0
\end{align*}
\end{teorema}

\setcounter{tcb@cnt@teorema}{19}
\begin{teorema}{Ausencia de Asociatividad}{no_asoc_deriv_eng}
\[ (a \uparrow b) \uparrow c \neq a \uparrow (b \uparrow c) \qquad \text{y} \qquad (a \downarrow b) \downarrow c \neq a \downarrow (b \downarrow c) \]
\end{teorema}

\setcounter{tcb@cnt@definicion}{6}
\begin{definicion}{NAND y NOR de 3 entradas}{n_entradas_eng}
\[ \uparrow(a,b,c) = \overline{a \cdot b \cdot c} \qquad \text{y} \qquad \downarrow(a,b,c) = \overline{a + b + c} \]
\end{definicion}

\setcounter{tcb@cnt@teorema}{20}
\begin{teorema}{NAND/NOR múltiple vs cascada binaria}{multiple_vs_binaria_eng}
\begin{align*}
    \uparrow(a,b,c) &\neq (a \uparrow b) \uparrow c \qquad & \uparrow(a,b,c) &\neq a \uparrow (b \uparrow c) \\
    \downarrow(a,b,c) &\neq (a \downarrow b) \downarrow c \qquad & \downarrow(a,b,c) &\neq a \downarrow (b \downarrow c)
\end{align*}
\end{teorema}
\subsection*{Comportamiento de los Operadores XOR y XNOR}

\setcounter{tcb@cnt@definicion}{4}
\begin{definicion}{Definición de XOR y XNOR}{def_xor_xnor_eng}
\[ a \oplus b = (a \cdot \overline{b}) + (\overline{a} \cdot b) \qquad \text{y} \qquad a \odot b = \overline{a \oplus b} = (a \cdot b) + (\overline{a} \cdot \overline{b}) \]
\end{definicion}

\setcounter{tcb@cnt@teorema}{22}
\begin{teorema}{Conmutatividad}{conmut_xor_eng}
\[ a \oplus b = b \oplus a \qquad \text{y} \qquad a \odot b = b \odot a \]
\end{teorema}

\setcounter{tcb@cnt@teorema}{23}
\begin{teorema}{Elementos Neutros e Inversores}{neutros_xor_eng}
\begin{align*}
    a \oplus 0 &= a \qquad & a \odot 1 &= a \\
    a \oplus 1 &= \overline{a} \qquad & a \odot 0 &= \overline{a}
\end{align*}
\end{teorema}

\setcounter{tcb@cnt@teorema}{24}
\begin{teorema}{Elemento Inverso de sí mismo (Grupo Abeliano)}{idemp_nula_eng}
\[ a \oplus a = 0 \qquad \text{y} \qquad a \odot a = 1 \]
\end{teorema}

\setcounter{tcb@cnt@teorema}{25}
\begin{teorema}{Propiedades de Negación}{neg_xor_eng}
\[ \overline{a \oplus b} = \overline{a} \oplus b = a \oplus \overline{b} = a \odot b \]
\[ \overline{a \odot b} = \overline{a} \odot b = a \odot \overline{b} = a \oplus b \]
\end{teorema}

\setcounter{tcb@cnt@teorema}{26}
\begin{teorema}{Asociatividad y Generalización}{asoc_gen_xor_eng}
Ambos operadores son asociativos:
\[ a \oplus (b \oplus c) = (a \oplus b) \oplus c \qquad \text{y} \qquad a \odot (b \odot c) = (a \odot b) \odot c \]
Lo cual permite su generalización a un número arbitrario $n$ de entradas:
\[ \bigoplus_{i=1}^{n} x_i = x_1 \oplus x_2 \oplus \dots \oplus x_n \qquad \text{y} \qquad \bigodot_{i=1}^{n} x_i = x_1 \odot x_2 \odot \dots \odot x_n \]
\end{teorema}

\setcounter{tcb@cnt@teorema}{28}
\begin{teorema}{Distributividad con el producto y la suma}{dist_xor_eng}
\[ a \cdot (b \oplus c) = (a \cdot b) \oplus (a \cdot c) \qquad \text{y} \qquad a + (b \odot c) = (a + b) \odot (a + c) \]
\end{teorema}

\subsection*{Tablas de Verdad Bivaluadas}

Resumen de las tablas de operación del álgebra de Boole para el caso de dos elementos ($B=\{0,1\}$), transcritas al lenguaje ingenieril:

\begin{table}[h!]
\centering
\renewcommand{\arraystretch}{1.2}
\begin{tabular}{c||c}
\hline
$a$ & $\overline{a}$ \\
\hline\hline
$0$ & $1$ \\
$1$ & $0$ \\
\hline
\end{tabular}
\qquad
\begin{tabular}{cc||c|c|c|c|c|c}
\hline
$a$ & $b$ & $a + b$ & $a \cdot b$ & $a \uparrow b$ & $a \downarrow b$ & $a \oplus b$ & $a \odot b$ \\
\hline\hline
$0$ & $0$ & $0$ & $0$ & $1$ & $1$ & $0$ & $1$ \\
$0$ & $1$ & $1$ & $0$ & $1$ & $0$ & $1$ & $0$ \\
$1$ & $0$ & $1$ & $0$ & $1$ & $0$ & $1$ & $0$ \\
$1$ & $1$ & $1$ & $1$ & $0$ & $0$ & $0$ & $1$ \\
\hline
\end{tabular}
\end{table}

\makeatother
\end{document}
